\documentclass[11pt]{article}

\usepackage[margin=1in]{geometry}
\usepackage{amsmath,amssymb,amsthm,bm}
\usepackage{booktabs}
\usepackage{graphicx}
\usepackage[hidelinks]{hyperref}
\usepackage{microtype}

\newcommand{\E}{\mathbb{E}}
\newcommand{\R}{\mathbb{R}}
\newcommand{\tr}{\operatorname{tr}}
\newcommand{\Var}{\operatorname{Var}}
\newcommand{\Cov}{\operatorname{Cov}}
\newcommand{\bhat}{\widehat{\bm\beta}}
\newcommand{\bstar}{\bm\beta^{\star}}
\newcommand{\that}{\widehat{\bm\theta}}
\newcommand{\tstar}{\bm\theta^{\star}}

\title{Prediction and Pixel-Space Control from Ridge Regression\\
in Linear Feature Spaces}
\author{Theory note, revised with deterministic-equivalent feature-space analysis}
\date{\today}

\begin{document}
\maketitle

\begin{abstract}
We study a noisy linear teacher and a ridge-regression student.  The student
may be trained after a linear feature map, but accentuation is performed by
back-propagating its gradient to the input pixels.  The feature map creates two
distinct geometries: the feature covariance
$C=A\Sigma A^{\top}$ controls natural-image prediction, whereas
$G=AA^{\top}$ controls pixel-space gradient motion.  This distinction gives
closed-form finite-sample expressions for generalization, own-path
accentuation, and peer review.  Weighted one- and two-resolvent deterministic
equivalents (DEs) give aggregate prediction and accentuation quantities for an
arbitrary fixed linear map.  When the map is spectral and aligned with the
eigenvectors of $\Sigma$, these results further resolve into per-mode weight
bias and variance.  The same DE risk can also approximate cross-validation and
select the ridge penalty.  A Gaussian distributional DE propagates fitted-weight
fluctuations through the nonlinear accentuation ratios when their quadratic
denominators do not concentrate.  Full PCA, whitening, and top-PC truncation
emerge as transparent special cases.
\end{abstract}

\section{Teacher, features, and ridge student}

Let the centered input and noisy response be
\begin{equation}
  \bm x\sim\mathcal N(0,\Sigma),\qquad
  y=(\bstar)^{\top}\bm x+\sigma\epsilon,
  \qquad \epsilon\sim\mathcal N(0,1).
\end{equation}
The noiseless teacher signal variance is
\begin{equation}
  S=(\bstar)^{\top}\Sigma\bstar.
\end{equation}
We first map the input to $p$ linear features,
\begin{equation}
  \bm z=A\bm x\in\R^p,
\end{equation}
and fit ridge regression in feature coordinates.  With the training feature
matrix $Z\in\R^{n\times p}$ and
$\widehat C=Z^{\top}Z/n$, the normalized-loss convention is
\begin{equation}
  \that_\lambda
  =\arg\min_{\bm\theta}
   \left\{\frac{1}{n}\lVert\bm y-Z\bm\theta\rVert^2
                 +\lambda\lVert\bm\theta\rVert^2\right\}
  =(\widehat C+\lambda I)^{-1}\frac{Z^{\top}\bm y}{n}.
  \label{eq:feature-ridge}
\end{equation}
The resulting model as a function of pixels is
\begin{equation}
  \widehat f(\bm x)=\that^{\top}A\bm x=\bhat^{\top}\bm x,
  \qquad \bhat=A^{\top}\that.
  \label{eq:pixel-weight}
\end{equation}
Thus feature-space estimation and pixel-space attribution are not the same
object.

\subsection{Best population predictor in the feature span}

Define
\begin{equation}
  C=A\Sigma A^{\top},\qquad
  \bm h=A\Sigma\bstar.
\end{equation}
The best linear feature coefficient and its unrepresented residual are
\begin{align}
  \tstar &= C^{\dagger}\bm h,\\
  r(\bm x)&=(\bstar)^{\top}\bm x-(\tstar)^{\top}A\bm x.
\end{align}
By the normal equations, $\Cov(\bm z,r)=0$.  Since the joint distribution is
Gaussian, the residual is independent of the retained features.  Its variance
is
\begin{equation}
  S_{\perp}=\Var(r)=S-(\tstar)^{\top}C\tstar.
  \label{eq:residual-signal}
\end{equation}
Consequently, truncated features reduce the training model to
\begin{equation}
  y=(\tstar)^{\top}\bm z+\xi,\qquad
  \Var(\xi)=\sigma_{\rm eff}^2=\sigma^2+S_{\perp}.
  \label{eq:effective-noise}
\end{equation}
The omitted teacher signal is simultaneously irreducible generalization error
and additional independent training noise.

It is important not to confuse this prediction-optimal coefficient with the
Euclidean projection of the teacher weight.  Define
\begin{equation}
  G=A A^{\top},\qquad \bm q=A\bstar,
\end{equation}
and let
\begin{equation}
  \bm\theta_{\rm E}=G^{\dagger}\bm q
\end{equation}
minimize $\lVert\bstar-A^{\top}\bm\theta\rVert^2$.  Then
\begin{align}
  A^{\top}\tstar
    &=A^{\top}(A\Sigma A^{\top})^{\dagger}A\Sigma\bstar,
    &&\text{prediction-optimal,}\label{eq:sigma-projection}\\
  A^{\top}\bm\theta_{\rm E}
    &=A^{\top}(AA^{\top})^{\dagger}A\bstar,
    &&\text{Euclidean-weight-optimal.}\label{eq:euclidean-projection}
\end{align}
The first is the $\Sigma$-orthogonal projection onto
$\operatorname{range}(A^{\top})$; the second is the ordinary Euclidean
projection.  They agree for isotropic inputs, for an invariant/aligned retained
subspace, or when this particular teacher has identical projections, but not
in general.

\section{Two geometries: prediction versus pixel control}

Natural-image prediction is governed by $C=A\Sigma A^{\top}$.  Pixel-space
gradient motion is governed by the distinct matrix
\begin{equation}
  G=AA^{\top},\qquad \bm q=A\bstar.
  \label{eq:control-geometry}
\end{equation}
The vector $\bm h=A\Sigma\bstar$ determines population prediction, whereas
$\bm q=A\bstar$ determines the teacher's response to a back-propagated pixel
gradient.  Confusing these two vectors hides the prediction--control
dissociation.

\subsection{Natural-distribution metrics}

For a fixed fitted model, the noiseless population prediction error is
\begin{equation}
  E_{\rm gen}(\that)
  =S_{\perp}+(\that-\tstar)^{\top}C(\that-\tstar).
  \label{eq:exact-gen-error}
\end{equation}
Equivalently, if
\begin{equation}
  N_{\rm gen}=\bm h^{\top}\that,
  \qquad D_{\rm gen}=\that^{\top}C\that,
\end{equation}
then
\begin{equation}
  \boxed{E_{\rm gen}=S-2N_{\rm gen}+D_{\rm gen}.}
  \label{eq:gen-quadratic-decomposition}
\end{equation}
This is an ordinary positive quadratic loss.  If performance is measured
against a newly observed noisy response rather than the noiseless teacher,
one simply adds the irreducible test noise $\sigma^2$.
We use the noiseless-teacher coefficient of determination
\begin{equation}
  R^2_{\rm gen}=1-\frac{E_{\rm gen}}{S}.
  \label{eq:r2-gen}
\end{equation}
Its denominator $S$ is a fixed population property.  Thus random fitted
quantities cannot make the normalization singular.  Generalization $R^2$ can
be negative when $E_{\rm gen}>S$, but it changes continuously with the fitted
weight and does not acquire the ratio instability of own-path accentuation.
The slope of true response on predicted response over natural inputs is
\begin{equation}
  b_{\rm gen}
  =\frac{\Cov((\bstar)^{\top}\bm x,\that^{\top}A\bm x)}
         {\Var(\that^{\top}A\bm x)}
  =\frac{\bm h^{\top}\that}{\that^{\top}C\that}.
  \label{eq:slope-gen-exact}
\end{equation}

\subsection{Own-path accentuation metrics}

Gradient accentuation from a seed $\bm x_0$ follows
\begin{equation}
  \bm x(t)=\bm x_0+t\nabla_{\bm x}\widehat f(\bm x)
           =\bm x_0+tA^{\top}\that.
  \label{eq:feature-accent-path}
\end{equation}
The student and teacher changes along this path are, respectively,
\begin{align}
  \Delta\widehat f(t)&=t\that^{\top}G\that,\\
  \Delta f^{\star}(t)&=t\bm q^{\top}\that.
  \label{eq:path-gains}
\end{align}
Therefore the exact true-on-predicted slope along the path is
\begin{equation}
  b_{\rm acc}(\that)
  =\frac{\bm q^{\top}\that}{\that^{\top}G\that}.
  \label{eq:slope-acc-exact}
\end{equation}
Suppose the experimenter chooses the distribution of $t$ so that the
student's pathwise response variance equals the natural teacher variance $S$:
\begin{equation}
  \Var(t)=\frac{S}{(\that^{\top}G\that)^2}.
\end{equation}
The accentuation calibration error is then
\begin{equation}
  \frac{E_{\rm acc}}{S}
  =\left(1-b_{\rm acc}\right)^2
  =\left(1-\frac{\bm q^{\top}\that}
                   {\that^{\top}G\that}\right)^2.
  \label{eq:acc-error-exact}
\end{equation}
When the denominator of $R^2$ is the variance of the \emph{true} response on
the generated path, the exact pathwise statistic becomes
\begin{equation}
  R^2_{\rm acc}
  =1-\left(1-\frac{\that^{\top}G\that}
                        {\bm q^{\top}\that}\right)^2.
  \label{eq:r2-acc-exact}
\end{equation}
Unlike $E_{\rm acc}/S$, this statistic is singular when the true path gain
$\bm q^{\top}\that$ is close to zero.  Hence it can be extremely negative even
when other summary metrics appear bounded.

\subsection{Two distinct accentuation breakdowns}

For clarity, write
\begin{equation}
  N_{\rm acc}=\bm q^{\top}\that,
  \qquad D_{\rm acc}=\that^{\top}G\that\geq 0.
\end{equation}
At a fixed, unnormalized path distribution, the raw discrepancy is simply
\begin{equation}
  \E_t[(\Delta\widehat f-\Delta f^{\star})^2]
  =\E[t^2](D_{\rm acc}-N_{\rm acc})^2,
  \label{eq:raw-path-loss}
\end{equation}
which is again a regular quadratic quantity.  The dramatic behavior enters
when the path is normalized by one of its fitted response variances:
\begin{align}
  \frac{E_{\rm acc}}{S}
    &=\left(1-\frac{N_{\rm acc}}{D_{\rm acc}}\right)^2,
    &\text{singular when }D_{\rm acc}\to0,\label{eq:eacc-singularity}\\
  R^2_{\rm acc}
    &=1-\left(1-\frac{D_{\rm acc}}{N_{\rm acc}}\right)^2,
    &\text{singular when }N_{\rm acc}\to0.\label{eq:r2acc-singularity}
\end{align}
These are opposite singular limits.  If $N_{\rm acc}\to0$ while
$D_{\rm acc}$ remains appreciable, then
$R^2_{\rm acc}\sim-(D_{\rm acc}/N_{\rm acc})^2\to-\infty$, whereas
$E_{\rm acc}/S\to1$.  If $D_{\rm acc}\to0$ while $N_{\rm acc}$ decays more
slowly, the student-calibrated loss diverges, whereas $R^2_{\rm acc}\to0$.
Moreover, $N_{\rm acc}<0$ means that the teacher moves opposite to the
student along the generated path; then $R^2_{\rm acc}<0$ automatically.

There are two mechanisms for approaching these bad regimes.  First, finite
samples, response noise, and cross-validated regularization can shrink or
randomize the signed linear form $N_{\rm acc}$.  At the same time, estimation
variance adds a positive contribution to $D_{\rm acc}$, so the ratio
$D_{\rm acc}/N_{\rm acc}$ can become very large.  Second, a non-aligned
dimension-reducing feature map creates a structural mismatch even with
infinite noiseless data.  In the limit $n\to\infty$, $\lambda\to0$,
\begin{equation}
  \that\longrightarrow\tstar=C^{\dagger}\bm h,
\end{equation}
and hence
\begin{equation}
  N_{\rm acc,\infty}=\bm q^{\top}\tstar,
  \qquad
  D_{\rm acc,\infty}=(\tstar)^{\top}G\tstar.
  \label{eq:structural-acc-gains}
\end{equation}
Because $\bm q=G\bm\theta_{\rm E}$ on the effective feature support,
\begin{equation}
  N_{\rm acc,\infty}-D_{\rm acc,\infty}
  =(\tstar)^{\top}G(\bm\theta_{\rm E}-\tstar).
  \label{eq:projection-mismatch-control}
\end{equation}
Prediction training selects $\tstar$, whereas exact pixel-gradient control
would prefer $\bm\theta_{\rm E}$.  The two need not agree, and the signed
quantity $\bm q^{\top}\tstar$ can be small or even negative despite small
natural generalization error.

This structural mechanism is absent in direct pixel regression.  With $A=I$,
\begin{equation}
  C=\Sigma,\quad G=I,\quad \bm h=\Sigma\bstar,
  \quad\bm q=\bstar,\quad S_{\perp}=0,
  \quad\tstar=\bm\theta_{\rm E}=\bstar.
\end{equation}
Thus infinite-data, zero-ridge pixel regression has
$N_{\rm acc}=D_{\rm acc}=\lVert\bstar\rVert^2$ and perfect own-path
calibration.  Its breakdown is entirely a finite-sample/noise/regularization
effect, often driven by low-variance input directions.  By contrast, a fixed
non-aligned truncated feature map can have both that statistical effect and
an irreducible projection mismatch.  An aligned PCA truncation is an
important intermediate case: it has $S_{\perp}>0$ in general, but its
prediction- and Euclidean-optimal weights agree on the retained modes, so its
infinite-data own-path calibration remains exact.

\section{Deterministic equivalents for an arbitrary fixed feature map}
\label{sec:fixed-feature-de}

The preceding exact formulas do not require $C$ and $G$ to commute.  Their
leading aggregate DEs also remain tractable for a deterministic, non-aligned
map.  Assume $p/n$ has a finite positive limit and the relevant operator norms
and deterministic quadratic forms remain bounded.  Let $\kappa$ solve
\begin{equation}
  \kappa-\lambda
  =\frac{\kappa}{n}\tr[C(C+\kappa I)^{-1}],
  \label{eq:general-kappa-fixed-point}
\end{equation}
and define
\begin{equation}
  \mathcal R=(C+\kappa I)^{-1},\qquad
  H=C\mathcal R,\qquad
  \bm\mu=H\tstar,
\end{equation}
as well as
\begin{equation}
  v_0=
  \frac{\kappa^2(\tstar)^{\top}\mathcal R C\mathcal R\tstar
        +\sigma_{\rm eff}^2}
       {n-\tr(C^2\mathcal R^2)},
  \qquad \sigma_{\rm eff}^2=\sigma^2+S_{\perp}.
  \label{eq:general-v0}
\end{equation}
Then the ordinary generalization quantities have DEs
\begin{align}
  N_{\rm gen}^{\rm DE}
    &=\bm h^{\top}\bm\mu,\\
  D_{\rm gen}^{\rm DE}
    &=\bm\mu^{\top}C\bm\mu
      +v_0\tr(C^2\mathcal R^2),\\
  E_{\rm gen}^{\rm DE}
    &=S-2N_{\rm gen}^{\rm DE}+D_{\rm gen}^{\rm DE}\nonumber\\
    &=S_{\perp}+(\bm\mu-\tstar)^{\top}C(\bm\mu-\tstar)
      +v_0\tr(C^2\mathcal R^2).
  \label{eq:general-feature-gen-de}
\end{align}
The first two terms in the final line are approximation bias and omitted
signal; the last is estimation variance.

For pixel accentuation,
\begin{align}
  N_{\rm acc}^{\rm DE}
    &=\bm q^{\top}\bm\mu,\\
  D_{\rm acc}^{\rm DE}
    &=\bm\mu^{\top}G\bm\mu
      +v_0\tr(G\mathcal R C\mathcal R).
  \label{eq:general-feature-acc-de}
\end{align}
The second line follows from the weighted two-resolvent identity
\begin{equation}
  Q_{\lambda}GQ_{\lambda}\ \asymp\
  \left(\frac{\kappa}{\lambda}\right)^2
  \left[
    \mathcal R G\mathcal R
    +\frac{n^{-1}\tr(G\mathcal R C\mathcal R)}
           {1-n^{-1}\tr(C^2\mathcal R^2)}
      \mathcal R C\mathcal R
  \right],
  \label{eq:weighted-two-resolvent}
\end{equation}
where $Q_{\lambda}=(\widehat C+\lambda I)^{-1}$.  The correction after
$\mathcal R G\mathcal R$ is essential: replacing the two random resolvents
independently by their one-point DE misses the random-design contribution.
Equations~\eqref{eq:general-feature-gen-de} and
\eqref{eq:general-feature-acc-de} reduce to the componentwise formulas below
when $C$ and $G$ share an eigenbasis.

\section{Geometry of protected prediction and catastrophic accentuation}
\label{sec:protected-prediction-bad-control}

\subsection{Exact two-center geometry}

The cleanest geometric description uses the two feature-space inner products
\begin{equation}
  \langle\bm u,\bm v\rangle_C=\bm u^{\top}C\bm v,
  \qquad
  \langle\bm u,\bm v\rangle_G=\bm u^{\top}G\bm v.
\end{equation}
Natural prediction compares $\that$ with the prediction-optimal center
$\bm\theta_{\rm P}=\tstar=C^{\dagger}\bm h$ in the $C$ metric:
\begin{equation}
  E_{\rm gen}=S_{\perp}
  +\lVert\that-\bm\theta_{\rm P}\rVert_C^2.
  \label{eq:prediction-ellipsoid}
\end{equation}
Pixel control instead compares $\that$ with the Euclidean-weight-optimal
center $\bm\theta_{\rm E}=G^{\dagger}\bm q$ in the $G$ metric.  On the
effective feature support, $\bm q=G\bm\theta_{\rm E}$ and therefore
\begin{equation}
  N_{\rm acc}=\langle\that,\bm\theta_{\rm E}\rangle_G,
  \qquad
  D_{\rm acc}=\lVert\that\rVert_G^2.
\end{equation}

Let $B_{\rm E}=\lVert\bm\theta_{\rm E}\rVert_G^2>0$ and make the exact
$G$-orthogonal decomposition
\begin{equation}
  \that=a\bm\theta_{\rm E}+\bm u,
  \qquad
  \langle\bm u,\bm\theta_{\rm E}\rangle_G=0.
\end{equation}
Then
\begin{align}
  N_{\rm acc}&=aB_{\rm E},\\
  D_{\rm acc}&=a^2B_{\rm E}+\lVert\bm u\rVert_G^2,\\
  \frac{D_{\rm acc}}{N_{\rm acc}}
    &=a+\frac{\lVert\bm u\rVert_G^2}{aB_{\rm E}}.
  \label{eq:exact-control-angle-norm}
\end{align}
Equation~\eqref{eq:exact-control-angle-norm} separates two control failures.
The coefficient $a$ measures signed calibration along the realizable teacher;
$\bm u$ is pixel-gradient energy orthogonal to that teacher.  For $a>0$,
$R^2_{\rm acc}<0$ exactly when
\begin{equation}
  a+\frac{\lVert\bm u\rVert_G^2}{aB_{\rm E}}>2.
\end{equation}
It becomes arbitrarily negative when $a\to0$ with nonzero orthogonal energy,
or when the orthogonal energy becomes arbitrarily large relative to
$aB_{\rm E}$.  Natural prediction can nevertheless remain accurate because
Eq.~\eqref{eq:prediction-ellipsoid} measures displacement from a different
center, in a different metric.

\subsection{Sharp prediction cost of collapsing the true path gain}

Assume for the moment that $C$ is positive definite on the effective feature
support.  The smallest excess prediction loss among all fitted weights whose
true accentuation gain vanishes is
\begin{align}
  \epsilon_{N=0}
  &\equiv
  \min_{\bm\theta:\,\bm q^{\top}\bm\theta=0}
       \lVert\bm\theta-\bm\theta_{\rm P}\rVert_C^2\nonumber\\
  &=\boxed{
    \frac{(\bm q^{\top}\bm\theta_{\rm P})^2}
         {\bm q^{\top}C^{-1}\bm q}}.
  \label{eq:minimum-risk-to-zero-gain}
\end{align}
The minimizing fitted weight is
\begin{equation}
  \bm\theta_{N=0}
  =\bm\theta_{\rm P}
   -\frac{\bm q^{\top}\bm\theta_{\rm P}}
          {\bm q^{\top}C^{-1}\bm q}C^{-1}\bm q.
\end{equation}
Hence a nearly perfect natural predictor can sit arbitrarily close to the
$R^2_{\rm acc}$ singularity when both
\begin{equation}
  \frac{S_{\perp}}{S}\ll1,
  \qquad
  \frac{\epsilon_{N=0}}{S}\ll1.
  \label{eq:protected-prediction-singular-control}
\end{equation}
The denominator $\bm q^{\top}C^{-1}\bm q$ is the susceptibility of the true
path gain to a unit of prediction loss.  It is large when the control
functional $\bm q$ points strongly into low-variance feature directions.  For
singular $C$, the same statement holds after restricting to its support and
using the corresponding pseudoinverse, provided the constraint is reachable.

\subsection{Generalized low-variance amplification}

For any feature-space error $\bm e$, define the generalized Rayleigh quotient
\begin{equation}
  \rho(\bm e)
  =\frac{\bm e^{\top}G\bm e}{\bm e^{\top}C\bm e}
  =\frac{\lVert A^{\top}\bm e\rVert^2}
         {(A^{\top}\bm e)^{\top}\Sigma(A^{\top}\bm e)}.
  \label{eq:generalized-control-rayleigh}
\end{equation}
The most dangerous direction is measured by
\begin{equation}
  \rho_{\max}
  =\lambda_{\max}(C^{-1/2}GC^{-1/2}),
  \qquad G\preceq\rho_{\max}C.
  \label{eq:max-control-amplification}
\end{equation}
Thus
\begin{equation}
  \lVert\bm e\rVert_G^2
  \leq\rho_{\max}\lVert\bm e\rVert_C^2.
\end{equation}
When $\rho_{\max}$ is modest, good prediction controls pixel-gradient error.
When it is very large, a direction can be almost invisible to natural risk
but enormous in pixel space.  A useful sufficient protection diagnostic is
\begin{equation}
  \Delta_G
  =\lVert\bm\theta_{\rm P}-\bm\theta_{\rm E}\rVert_G
   +\sqrt{\rho_{\max}(E_{\rm gen}-S_{\perp})}.
  \label{eq:control-protection-radius}
\end{equation}
If $\Delta_G\ll\lVert\bm\theta_{\rm E}\rVert_G$, then prediction accuracy
forces $\that$ into a small $G$-ball around the control-optimal weight and
accentuation is protected.  If $\Delta_G$ is comparable to or larger than
$\lVert\bm\theta_{\rm E}\rVert_G$, natural prediction alone supplies no such
guarantee.  The first term is structural feature mismatch; the second is
statistical error amplified by low natural variance.

The DE has a scale-specific version of the same amplification.  Estimation
variance contributes
\begin{equation}
  V_{\rm gen}=v_0\tr(C^2\mathcal R^2),
  \qquad
  V_{\rm acc}=v_0\tr(G\mathcal R C\mathcal R),
\end{equation}
so their ratio is
\begin{equation}
  \mathcal A_{G/C}(\kappa)
  =\frac{\tr(G\mathcal R C\mathcal R)}
         {\tr(C^2\mathcal R^2)}.
  \label{eq:de-control-amplification}
\end{equation}
Prediction can remain protected while accentuation variance is large when
$V_{\rm gen}\ll S$ but
$V_{\rm acc}$ is comparable to $|N_{\rm acc}|$; a large
$\mathcal A_{G/C}(\kappa)$ makes these inequalities compatible.

\subsection{An exact DE identity linking generalization risk to control
variance}

The same scalar $v_0$ that controls fitted-weight fluctuations has a direct
interpretation in terms of generalization error.  Since
\begin{equation}
  \bm\mu-\tstar=-\kappa\mathcal R\tstar,
\end{equation}
Eq.~\eqref{eq:general-feature-gen-de} gives
\begin{equation}
  E_{\rm gen}^{\rm DE}+\sigma^2
  =S_{\perp}+\sigma^2
   +\kappa^2(\tstar)^{\top}\mathcal R C\mathcal R\tstar
   +v_0\tr(C^2\mathcal R^2).
\end{equation}
Using the definition of $v_0$ in Eq.~\eqref{eq:general-v0}, the right side
simplifies exactly to
\begin{equation}
  \boxed{E_{\rm gen}^{\rm DE}+\sigma^2=n v_0.}
  \label{eq:egen-v0-identity}
\end{equation}
Here $E_{\rm gen}$ is error against the noiseless teacher and already includes
$S_{\perp}$; consequently only the observation noise $\sigma^2$, rather than
$\sigma_{\rm eff}^2$, is added on the left.  The sum
$E_{\rm gen}^{\rm DE}+\sigma^2$ is the DE test risk against a newly observed
noisy response.

Substituting Eq.~\eqref{eq:egen-v0-identity} into the second moments yields
\begin{align}
  D_{\rm gen}^{\rm DE}
    &=\bm\mu^{\top}C\bm\mu
      +\frac{E_{\rm gen}^{\rm DE}+\sigma^2}{n}
       \tr(C^2\mathcal R^2),\\
  \boxed{
  D_{\rm acc}^{\rm DE}
    =\bm\mu^{\top}G\bm\mu
      +\frac{E_{\rm gen}^{\rm DE}+\sigma^2}{n}
       \tr(G\mathcal R C\mathcal R).}
  \label{eq:dacc-egen-identity}
\end{align}
Thus generalization risk fixes the global fluctuation amplitude, but the
evaluation geometry decides how strongly those fluctuations enter a given
quadratic statistic.

For direct pixel regression, set $C=\Sigma$, $G=I$, and
$\tstar=\bstar$.  With
\begin{align}
  B_{a,b}
    &=(\bstar)^{\top}\Sigma^a(\Sigma+\kappa I)^{-b}\bstar,\\
  \mathrm{df}_{a,b}
    &=\tr[\Sigma^a(\Sigma+\kappa I)^{-b}],
\end{align}
Eq.~\eqref{eq:dacc-egen-identity} becomes
\begin{equation}
  \boxed{
  \mu_D^{\rm DE}
    =B_{2,2}
     +\frac{\mathrm{df}_{1,2}}{n}
       (E_{\rm gen}^{\rm DE}+\sigma^2).}
  \label{eq:pixel-d-egen-identity}
\end{equation}
For a feature model, this same unweighted expression is valid only if
$G=I_p$ or if $D$ denotes feature-coefficient norm.  Pixel-space
back-propagation requires the $G$-weighted form in
Eq.~\eqref{eq:dacc-egen-identity}; for commuting $C$ and $G$ it reads
\begin{equation}
  D_{\rm acc}^{\rm DE}
    =B_{2,2}^{G}
     +\frac{d_{1,2}^{G}}{n}
       (E_{\rm gen}^{\rm DE}+\sigma^2).
  \label{eq:aligned-feature-d-egen-identity}
\end{equation}

Combining this identity with $N_{\rm acc}^{\rm DE}=\bm q^{\top}\bm\mu$
gives the useful decomposition
\begin{equation}
  \boxed{
  z_{\rm acc}^{\rm DE}
  =\frac{\bm\mu^{\top}G\bm\mu-\bm q^{\top}\bm\mu}
         {\bm q^{\top}\bm\mu}
   +\frac{E_{\rm gen}^{\rm DE}+\sigma^2}{n}
    \frac{\tr(G\mathcal R C\mathcal R)}
         {\bm q^{\top}\bm\mu}.}
  \label{eq:z-egen-decomposition}
\end{equation}
The first term is structural mismatch plus shrinkage; the second is
generalization-linked variance inflation.  Therefore small
$E_{\rm gen}^{\rm DE}$ does not by itself protect accentuation: the second
term can still be large if the weighted trace is large or the signed true
path gain $\bm q^{\top}\bm\mu$ is small.  Conversely, if the first term is
small and the multiplier
\begin{equation}
  \frac{\tr(G\mathcal R C\mathcal R)}
       {n|\bm q^{\top}\bm\mu|}
\end{equation}
is controlled, low generalization risk does protect own-path calibration.
For positive $N_{\rm acc}^{\rm DE}$, the leading DE predicts negative
$R^2_{\rm acc}$ exactly when
\begin{equation}
  \bm\mu^{\top}G\bm\mu
  +\frac{E_{\rm gen}^{\rm DE}+\sigma^2}{n}
       \tr(G\mathcal R C\mathcal R)
  >2\bm q^{\top}\bm\mu.
  \label{eq:negative-r2-threshold-egen}
\end{equation}

\subsection{CV-selected simplification of the accentuation ratio}

Let
\begin{equation}
  z_{\rm acc}^{\rm DE}
  =\frac{D_{\rm acc}^{\rm DE}}{N_{\rm acc}^{\rm DE}}-1.
\end{equation}
For an arbitrary fixed map, Eqs.~\eqref{eq:general-feature-acc-de} give
\begin{equation}
  z_{\rm acc}^{\rm DE}
  =\frac{\bm\mu^{\top}G\bm\mu-\bm q^{\top}\bm\mu
          +v_0T_G}
         {\bm q^{\top}\bm\mu},
  \qquad
  T_G=\tr(G\mathcal R C\mathcal R).
  \label{eq:general-feature-z}
\end{equation}
At an interior minimizer of the DE generalization risk, define
\begin{equation}
  B_{2,3}^{C}=(\tstar)^{\top}C^2\mathcal R^3\tstar,
  \qquad
  d_{2,3}^{C}=\tr(C^2\mathcal R^3).
\end{equation}
The first-order stationarity condition implies
\begin{equation}
  v_0=\kappa\frac{B_{2,3}^{C}}{d_{2,3}^{C}}.
  \label{eq:cv-v0-stationarity}
\end{equation}
Using $\bm q=G\bm\theta_{\rm E}$, the CV-selected ratio therefore becomes
\begin{equation}
  \boxed{
  z_{\rm acc,CV}^{\rm DE}
  =\frac{\bm\mu^{\top}G(\bm\mu-\bm\theta_{\rm E})}
         {\bm q^{\top}\bm\mu}
   +\frac{\kappa T_G}{\bm q^{\top}\bm\mu}
      \frac{B_{2,3}^{C}}{d_{2,3}^{C}}.}
  \label{eq:general-feature-z-cv}
\end{equation}
The first term is deterministic shrinkage plus structural projection mismatch;
the second is variance inflation at the penalty selected to protect natural
prediction.  Cross-validation controls the latter only through $C$ and does
not optimize the $G$ geometry or remove the first term.

There is a closer analogue of the pixel-regression simplification when $C$
and $G$ commute and $\bm\theta_{\rm P}=\bm\theta_{\rm E}=\tstar$, as for an
aligned spectral feature map.  Define
\begin{align}
  B_{a,b}^{G}
    &=(\tstar)^{\top}G C^a(C+\kappa I)^{-b}\tstar,\\
  d_{a,b}^{G}
    &=\tr[G C^a(C+\kappa I)^{-b}].
\end{align}
Then
\begin{equation}
  N_{\rm acc}^{\rm DE}=B_{1,1}^{G},
  \qquad
  D_{\rm acc}^{\rm DE}=B_{2,2}^{G}+v_0d_{1,2}^{G},
\end{equation}
and the identity $B_{2,2}^{G}-B_{1,1}^{G}=-\kappa B_{1,2}^{G}$ yields
\begin{equation}
  \boxed{
  z_{\rm acc,CV}^{\rm DE}
  =\kappa\frac{d_{1,2}^{G}}{B_{1,1}^{G}}
   \left(
     \frac{B_{2,3}^{C}}{d_{2,3}^{C}}
     -\frac{B_{1,2}^{G}}{d_{1,2}^{G}}
   \right).}
  \label{eq:aligned-feature-z-cv}
\end{equation}
This has the same difference-of-weighted-averages structure as direct pixel
regression, but the CV average is selected in the natural prediction geometry
$C$, while the second average is weighted by pixel-control geometry $G$.
Setting $A=I$, so that $C=\Sigma$, $G=I$, and $\tstar=\bstar$, recovers the
pixel-space formula.

Finally, $R_{\rm acc}^{2,\rm DE}=1-(z_{\rm acc}^{\rm DE})^2$ is a typical-fit
plug-in, not generally the Monte Carlo mean.  The replacement of a random
ratio by a ratio of deterministic means requires relative concentration of
its denominator.  Precisely in the catastrophic regime where
$N_{\rm acc}$ approaches or crosses zero, that expansion is nonuniform and
rare realizations can dominate the empirical mean.  Medians, quantiles, and a
joint fluctuation theory for $(N_{\rm acc},D_{\rm acc})$ are then essential.

\subsection{Gaussian distributional DE for nonlinear control metrics}
\label{sec:gaussian-distributional-de}

The leading DE determines the first moments of $N_{\rm acc}$ and
$D_{\rm acc}$, but own-path statistics are nonlinear functions of these two
random quantities.  A useful intermediate theory keeps the DE mean and
covariance of the fitted feature weight rather than immediately substituting
their aggregate expectations.  Write
\begin{equation}
  \that\ \dot\sim\ \mathcal N(\bm\mu,V),
  \qquad
  V=v_0\mathcal R C\mathcal R,
  \label{eq:gaussian-weight-de}
\end{equation}
where $\bm\mu$, $v_0$, and $\mathcal R$ are defined in
Section~\ref{sec:fixed-feature-de}.  This Gaussian approximation gives the
joint linear--quadratic moments
\begin{align}
  \mu_N&=\E[N_{\rm acc}]=\bm q^{\top}\bm\mu,
  &\Var(N_{\rm acc})&=\bm q^{\top}V\bm q,\\
  \mu_D&=\E[D_{\rm acc}]
       =\bm\mu^{\top}G\bm\mu+\tr(GV),
  &\Var(D_{\rm acc})
       &=2\tr[(GV)^2]+4\bm\mu^{\top}GVG\bm\mu,\\
  &&\Cov(N_{\rm acc},D_{\rm acc})
       &=2\bm q^{\top}VG\bm\mu.
  \label{eq:gaussian-nd-moments}
\end{align}
The leading formulas in Eq.~\eqref{eq:general-feature-acc-de} use
$\mu_N/\mu_D$ and $\mu_D/\mu_N$.  They approximate per-fit nonlinear
statistics only when the relevant random denominators concentrate.  Two
direct diagnostics are
\begin{equation}
  \mathrm{CV}_D
  =\frac{\sqrt{\Var(D_{\rm acc})}}{\mu_D},
  \qquad
  \mathrm{CV}_N
  =\frac{\sqrt{\Var(N_{\rm acc})}}{|\mu_N|}.
  \label{eq:nd-concentration-diagnostics}
\end{equation}
For example, a second-order delta expansion gives
\begin{equation}
  \E\left[\frac{N_{\rm acc}}{D_{\rm acc}}\right]
  \approx
  \frac{\mu_N}{\mu_D}
  -\frac{\Cov(N_{\rm acc},D_{\rm acc})}{\mu_D^2}
  +\frac{\mu_N\Var(D_{\rm acc})}{\mu_D^3}.
  \label{eq:delta-slope-acc}
\end{equation}
This is informative for small $\mathrm{CV}_D$, but is not a controlled
expansion when $\mathrm{CV}_D=O(1)$.  In that regime one should instead draw
\begin{equation}
  \that^{(b)}=\bm\mu+L\bm z^{(b)},
  \qquad LL^{\top}=V,
  \qquad \bm z^{(b)}\sim\mathcal N(0,I_p),
  \label{eq:gaussian-de-sampling}
\end{equation}
and propagate every draw through the exact per-fit definitions
\begin{align}
  b_{\rm acc}^{(b)}
    &=\frac{\bm q^{\top}\that^{(b)}}
            {(\that^{(b)})^{\top}G\that^{(b)}},\\
  \frac{E_{\rm acc}^{(b)}}{S}
    &=(1-b_{\rm acc}^{(b)})^2,\\
  R_{\rm acc}^{2,(b)}
    &=1-\left(
      1-\frac{(\that^{(b)})^{\top}G\that^{(b)}}
               {\bm q^{\top}\that^{(b)}}\right)^2.
  \label{eq:gaussian-de-control-metrics}
\end{align}
Means, medians, intervals, and tail probabilities of
Eq.~\eqref{eq:gaussian-de-control-metrics} define the \emph{Gaussian
distributional DE}.  It is not primarily a new correction to $\mu_N$ or
$\mu_D$; it retains the already-derived fitted-weight fluctuations while
respecting the nonlinear order of operations.

A useful summary of quadratic-form concentration is the effective rank
\begin{equation}
  r_{\rm eff}(V^{1/2}GV^{1/2})
  =\frac{\tr(GV)^2}{\tr[(GV)^2]}.
  \label{eq:control-effective-rank}
\end{equation}
When the random part dominates $D_{\rm acc}$,
$\mathrm{CV}_D\simeq\sqrt{2/r_{\rm eff}}$.  Thus feature dimension $p$ is
not itself the concentration parameter.  Increasing $p$ helps only if it adds
enough comparable control directions to increase effective rank; one enormous
low-variance direction can keep $r_{\rm eff}$ near one even in a nominally
wide feature space.

\paragraph{Van Hateren dimension-scaling check.}
We tested the same prediction-equivalent Givens rotation at
$p\in\{16,128,500,1000\}$ with $n=1000$.  The original $p=16$ map was kept
exactly, and larger maps appended disjoint directions orthogonal to both
$\Sigma^{1/2}\bstar$ and $\Sigma^{-1/2}\bstar$.  Hence $C$, the realizable
teacher, and the fixed DE-CV prediction problem were invariant along every
rotation.  At the low-variance endpoint and $\sigma^2/S=0.01$, the effective
rank increased from $1.08$ to $34.4$, $272$, and $657$, while
$\mathrm{CV}_D$ decreased from $1.20$ to $0.237$, $0.086$, and $0.055$.
The relative leading-DE discrepancies for the accentuation slope were
$60.7\%$, $4.6\%$, $0.9\%$, and $0.5\%$; for median
$R^2_{\rm acc}$ they were $342\%$, $10.3\%$, $1.0\%$, and $0.1\%$.
The Gaussian distributional DE matched the exact finite-$n$ ridge Monte Carlo
throughout the low-effective-rank regime.

Crucially, concentration restored the \emph{accuracy of the leading DE}, not
the quality of control.  As $p$ grew, many prediction-null directions added
pixel-gradient energy: the accentuation slope became smaller and
$R^2_{\rm acc}$ became more negative even though the leading DE became nearly
exact.  This separates two questions that should not be conflated:
\begin{enumerate}
  \item Is the leading ratio-of-moments DE statistically sufficient?  This is
        controlled by $\mathrm{CV}_D$, $\mathrm{CV}_N$, and effective rank.
  \item Is accentuation well calibrated?  This is controlled by the relative
        geometry of $C$, $G$, $\bm h$, and $\bm q$.
\end{enumerate}

\begin{figure}[t]
  \centering
  \includegraphics[width=0.92\textwidth]{figures/null_rotations/vanhateren_null_leading_de_adequacy.png}
  \caption{Low-variance endpoint of the Van Hateren prediction-null rotation.
  As feature dimension adds comparable control directions, the effective rank
  rises and the Gaussian distributional DE, leading DE, and exact finite-$n$
  ridge Monte Carlo converge.  Simultaneously, the control slope approaches
  zero and $R^2_{\rm acc}$ becomes increasingly negative.}
  \label{fig:null-dimension-leading-de}
\end{figure}

\section{Spectral feature maps and deterministic equivalents}
\label{sec:spectral-de}

We now specialize to the setting implemented in the accompanying numerical
experiments.  Write the input eigendecomposition as
\begin{equation}
  \Sigma=U\,\mathrm{diag}(s_1,\ldots,s_d)U^{\top},
  \qquad b_j=\bm u_j^{\top}\bstar.
\end{equation}
For an index set $\mathcal J$, let the feature map retain and rescale PCs,
\begin{equation}
  z_j=a_j\bm u_j^{\top}\bm x,\qquad j\in\mathcal J.
\end{equation}
In the common eigenbasis of prediction and control, define
\begin{equation}
  c_j=a_j^2s_j,\qquad
  g_j=a_j^2,\qquad
  \theta_j^{\star}=\frac{b_j}{a_j}.
  \label{eq:spectral-feature-quantities}
\end{equation}
Then $C=\mathrm{diag}(c_j)$, $G=\mathrm{diag}(g_j)$, and
$q_j=g_j\theta_j^{\star}$.  The omitted signal and effective noise are
\begin{equation}
  S_{\perp}=\sum_{j\notin\mathcal J}s_jb_j^2,qquad
  \sigma_{\rm eff}^2=\sigma^2+S_{\perp}.
\end{equation}

We take $p,n\to\infty$ with $p/n\to\gamma$.  The effective ridge parameter
$\kappa=\kappa(\lambda)$ solves the Silverstein fixed-point equation
\begin{equation}
  \kappa-\lambda
  =\frac{\kappa}{n}\sum_{j\in\mathcal J}
       \frac{c_j}{c_j+\kappa}.
  \label{eq:kappa-fixed-point}
\end{equation}
Equivalently, the right side is
$\gamma\kappa\int c/(c+\kappa)\,d\nu(c)$ for the empirical feature spectrum
$\nu$.

Define
\begin{align}
  h_j&=\frac{c_j}{c_j+\kappa},\\
  \mathrm{df}_2(\kappa)&=\sum_j\frac{c_j^2}{(c_j+\kappa)^2},\\
  C_{\rm sig}(\kappa)&=
       \sum_j\frac{c_j(\theta_j^{\star})^2}{(c_j+\kappa)^2},\\
  v_0(\kappa)&=
       \frac{\kappa^2C_{\rm sig}(\kappa)+\sigma_{\rm eff}^2}
            {n-\mathrm{df}_2(\kappa)}.
  \label{eq:de-global-quantities}
\end{align}
The leading componentwise first and second moments of the fitted feature
weight are
\begin{align}
  m_j&\equiv\E[\widehat\theta_j]
       \simeq h_j\theta_j^{\star},\\
  v_j&\equiv\Var(\widehat\theta_j)
       \simeq v_0\frac{c_j}{(c_j+\kappa)^2},\\
  \E[\widehat\theta_j^2]&\simeq m_j^2+v_j.
  \label{eq:de-component-moments}
\end{align}
Thus the per-feature coefficient error is
\begin{equation}
  e_j\equiv\E[(\widehat\theta_j-\theta_j^{\star})^2]
  \simeq (1-h_j)^2(\theta_j^{\star})^2+v_j.
  \label{eq:de-per-feature-error}
\end{equation}
The first term is ridge shrinkage bias; the second combines finite-sample
signal and response-noise fluctuations.

\subsection{DE generalization metrics}

Aggregating the componentwise moments gives
\begin{align}
  E_{\rm gen}^{\rm DE}
    &=S_{\perp}+\sum_j c_je_j,
    \label{eq:de-gen-error}\\
  R_{\rm gen}^{2,\rm DE}
    &=1-\frac{E_{\rm gen}^{\rm DE}}{S},\\
  N_{\rm gen}^{\rm DE}
    &=\sum_j c_jm_j\theta_j^{\star},\\
  D_{\rm gen}^{\rm DE}
    &=\sum_j c_j(m_j^2+v_j),\\
  b_{\rm gen}^{\rm DE}
    &=\frac{N_{\rm gen}^{\rm DE}}{D_{\rm gen}^{\rm DE}}.
  \label{eq:de-gen-slope}
\end{align}

\subsection{DE pixel-accentuation metrics}

The control geometry changes only the aggregation weights:
\begin{align}
  N_{\rm acc}^{\rm DE}
    &=\sum_j g_jm_j\theta_j^{\star},\\
  D_{\rm acc}^{\rm DE}
    &=\sum_j g_j(m_j^2+v_j),\\
  b_{\rm acc}^{\rm DE}
    &=\frac{N_{\rm acc}^{\rm DE}}{D_{\rm acc}^{\rm DE}},
    \label{eq:de-acc-slope}\\
  \frac{E_{\rm acc}^{\rm DE}}{S}
    &=\left(1-b_{\rm acc}^{\rm DE}\right)^2,\\
  R_{\rm acc}^{2,\rm DE}
    &=1-\left(1-\frac{D_{\rm acc}^{\rm DE}}
                            {N_{\rm acc}^{\rm DE}}\right)^2.
    \label{eq:de-acc-r2}
\end{align}
Equations~\eqref{eq:de-gen-error}--\eqref{eq:de-acc-r2} are the main
feature-space extension.  Natural generalization weights mode $j$ by $c_j$;
pixel accentuation weights it by $g_j$.  The two agree only in special cases.

The DE above substitutes ratios of deterministic first and second moments for
the nonlinear per-fit ratios.  It predicts a typical high-dimensional fit but
is not, in general, the exact Monte Carlo mean of $R^2_{\rm acc}$.  This
distinction matters whenever $N_{\rm acc}$ can approach zero: rare fits then
dominate the empirical mean, while the median can remain close to the leading
DE.

\subsection{Input-space weight recovery}

The recovered pixel coefficient in retained mode $j$ is
$\widehat b_j=a_j\widehat\theta_j$.  Therefore
\begin{align}
  \E[\widehat b_j]&\simeq a_jm_j,\\
  \Var(\widehat b_j)&\simeq g_jv_j,\\
  \E\lVert\bhat-\bstar\rVert^2
    &\simeq \sum_{j\in\mathcal J}g_je_j
          +\sum_{j\notin\mathcal J}b_j^2.
  \label{eq:de-pixel-weight-error}
\end{align}
This formula makes the high-frequency mechanism explicit.  A low-variance
direction can make a negligible contribution $s_j e_j$ to prediction error
but a large contribution $g_je_j$ to the back-propagated pixel gradient.
The DE predicts the mean and variance envelope across modes; it does not
predict the signs or texture of one particular random fitted realization.

\section{Cross-validated ridge selection from the DE}

Many software packages parameterize ridge as
\begin{equation}
  \lVert\bm y-Z\bm\theta\rVert^2+\alpha\lVert\bm\theta\rVert^2,
\end{equation}
so that $\lambda=\alpha/n$ in Eq.~\eqref{eq:feature-ridge}.  For analytic
leave-one-out cross-validation, a DE approximation follows the actual
training-size change:
\begin{enumerate}
  \item set $n_{\rm tr}=n-1$;
  \item for every candidate $\alpha$, solve
        Eq.~\eqref{eq:kappa-fixed-point} with
        $\lambda=\alpha/n_{\rm tr}$ and aspect ratio $p/n_{\rm tr}$;
  \item compute $E_{\rm gen}^{\rm DE}(\alpha;n_{\rm tr})$ using
        $\sigma_{\rm eff}^2$;
  \item select
        $\alpha_{\rm DE-CV}=\arg\min_{\alpha}E_{\rm gen}^{\rm DE}$;
  \item refit the final $n$-sample DE with
        $\lambda=\alpha_{\rm DE-CV}/n$.
\end{enumerate}
For $K$-fold CV, replace $n-1$ by the fold training size.  A dense log-spaced
$\alpha$ grid is important: a coarse grid creates artificial jumps in the
selected regularization and therefore jagged downstream control metrics.
The DE selects ridge for natural predictive risk, just as ordinary CV does;
it does not secretly optimize accentuation.

\section{PCA, whitening, and top-PC truncation}

The spectral formulas expose the behavior of the three feature spaces used in
the FFHQ disk-teacher experiment.

\begin{table}[h]
\centering
\small
\begin{tabular}{@{}lllll@{}}
\toprule
Feature map & $a_j$ & $c_j$ & $g_j$ & omitted signal \\
\midrule
Full PCA & $1$ & $s_j$ & $1$ & $0$ \\
Full whitening & $s_j^{-1/2}$ & $1$ & $s_j^{-1}$ & $0$ \\
Top $K$ PCs & $1$, $j\le K$ & $s_j$ & $1$ &
  $\sum_{j>K}s_jb_j^2$ \\
\bottomrule
\end{tabular}
\caption{Prediction and pixel-control quantities for aligned linear features.}
\label{tab:feature-cases}
\end{table}

\paragraph{Full PCA.}
An orthogonal change of basis leaves ridge and the pixel-gradient norm
unchanged.  It is the pixel-space baseline written in interpretable spectral
coordinates.

\paragraph{Full whitening.}
Regression sees an isotropic covariance, $c_j=1$, but back-propagation weights
each mode by $g_j=1/s_j$.  Low-variance, high-frequency PCs can therefore
dominate $D_{\rm acc}$ even though they do not dominate natural prediction.
The feature representation may look well conditioned while pixel control is
ill conditioned.

\paragraph{Top-PC truncation.}
Discarding low-variance PCs introduces irreducible natural error
$S_{\perp}$, but completely prevents those directions from entering the
pixel-gradient path.  When the teacher signal is concentrated in the leading
PCs, a modest predictive cost can buy a large control benefit.  The optimal
cutoff depends on noise: higher response noise generally favors a smaller,
more signal-focused feature space.

\begin{figure}[t]
  \centering
  \includegraphics[width=\textwidth]{figures/feature_space/ffhq_linear_feature_gen_acc_gap.png}
  \caption{Measured-spectrum FFHQ disk-teacher experiment.  Curves are the
  leading DE; points are Monte Carlo ridge fits.  Full PCA develops a
  generalization--accentuation gap, top-100-PC truncation suppresses the
  harmful control directions, and full whitening produces a near-zero true
  path gain and extremely negative $R^2_{\rm acc}$.  The lower horizontal
  axis is response-noise standard deviation; the upper axis is
  $\sigma^2/S$.}
  \label{fig:feature-gap}
\end{figure}

\section{Peer review across models and feature spaces}

Let an evaluator have pixel weight
$\widehat{\bm w}_e=A_e^{\top}\widehat{\bm\theta}_e$ and a generator have
$\widehat{\bm w}_g=A_g^{\top}\widehat{\bm\theta}_g$.  Along the generator's
path $\bm x(t)=\bm x_0+t\widehat{\bm w}_g$, the teacher and evaluator gains
are
\begin{equation}
  T_g=(\bstar)^{\top}\widehat{\bm w}_g,\qquad
  D_{e\leftarrow g}=\widehat{\bm w}_e^{\top}\widehat{\bm w}_g
  =\widehat{\bm\theta}_e^{\top}A_eA_g^{\top}
       \widehat{\bm\theta}_g.
\end{equation}
If the generator scales its path to its own natural-response target range,
the evaluator's normalized peer error, slope, and $R^2$ are
\begin{align}
  \frac{E_{e\leftarrow g}}{S}
    &=\left(\frac{D_{e\leftarrow g}-T_g}
                    {\widehat{\bm w}_g^{\top}\widehat{\bm w}_g}\right)^2,\\
  b_{e\leftarrow g}&=\frac{T_g}{D_{e\leftarrow g}},\\
  R^2_{e\leftarrow g}
    &=1-\left(1-\frac{D_{e\leftarrow g}}{T_g}\right)^2.
  \label{eq:peer-r2-feature}
\end{align}
For two independently trained models, a leading peer DE can replace each
weight by its deterministic mean in $D_{e\leftarrow g}$ and $T_g$.  Models
trained on the same or overlapping samples require a joint two-resolvent DE
to retain their correlated fluctuations.  Equation~\eqref{eq:peer-r2-feature}
also clarifies the observed asymmetry: a generator can move in directions
that remain legible to a good evaluator, while the reverse path can have
nearly zero teacher gain.

\section{Interpretation and limitations}

The main conclusions are:
\begin{enumerate}
  \item Good natural-image prediction does not imply good pixel-space
        control.  Prediction uses $C$; control uses $G$.
  \item Low-variance modes are cheap under natural risk but can be expensive
        under pixel-gradient accentuation, especially after whitening.
  \item Top-PC truncation can reverse the generalization--accentuation gap by
        removing directions that contribute little teacher signal but large
        control variance.
  \item The feature-space DE predicts mean per-mode bias and variance,
        aggregate natural risk, typical control calibration, and DE-CV ridge
        selection.  It does not specify the visual pattern of a single random
        fit.
  \item $R^2_{\rm acc}$ and peer $R^2$ are nonlinear ratios.  Near-zero true
        path gain creates heavy tails, so empirical medians/IQRs and
        means/standard errors should both be reported.  When the numerator or
        quadratic denominator does not concentrate, the Gaussian
        distributional DE in Section~\ref{sec:gaussian-distributional-de}
        should replace a direct ratio-of-moments plug-in.
\end{enumerate}

The aggregate first-order formulas in
Section~\ref{sec:fixed-feature-de} allow a fixed linear map with noncommuting
$C$ and $G$ through a weighted two-resolvent DE.  The mode-resolved bias and
variance formulas in Section~\ref{sec:spectral-de} additionally assume an
aligned spectral map.  Full fluctuations of nonlinear ratios, joint
peer-review statistics, feature maps learned from the same training sample,
and nonlinear features still require joint multi-resolvent or operator-valued
deterministic equivalents.

A detailed derivation of the current noise-conditional
$\operatorname{Var}(R)$ correction, together with the joint-resolvent
covariance calculation needed for the full random-design variance, is given in
\path{notes/accentuation_varR_theory.tex}.

\end{document}
